\documentclass[11pt]{article}
\usepackage{geometry}
\geometry{margin=1in}

\title{\textbf{MODEL README}}
\author{Mitchell Ballif and Ricardo Reyes-Heroles}
\date{October 2026}

\begin{document}

\maketitle



\section*{Programs}

\begin{itemize}

\item \texttt{run.m} --- solves the terminal steady state and the
transition from 1965 onto it, once at the recovered wedges and once with a
wedge frozen, then prints the counterfactual. Section 1, Steps 1 and 2.

\item \texttt{variables.m} --- names the two regions and three sectors and
fixes their index positions.

\item \texttt{parameters.m} --- the parameters not recovered from the data,
and the solver tolerances.

\item \texttt{loaddata.m} --- reads wedges from  processed panels, 1965 to 2014. 

\item \texttt{padyears.m} --- extends the sample to the steady state year

\item \texttt{baseline.m} ---Baseline levels from value added data,
consumption expenditure, and the final demand shares of Section 3.

\item \texttt{wedges.m} --- Returns wedge data or create new frozen wedge path

\item \texttt{static.m} --- Given a guess in wage,  static solves equilibrium conditions for a given year
\item \texttt{steadystate.m} --- the terminal steady state, by two nested
loops over capital and wages. Section 2.

\item \texttt{dynamic.m} --- the equilibrium transition. Guesses capital and
consumption expenditure, clears each period with \texttt{static.m}, 

\item \texttt{makefigures.m} --- the figures, from
\texttt{results/solution.mat}.

\item \texttt{checknote.m} --- re-solves the path and checks it against all 26
conditions in \texttt{EQconditions.tex}.

\end{itemize}

\section*{Other files}

\begin{itemize}

\item \texttt{EQconditions.tex} --- the model. Section 1 is the transition
algorithm, Section 2 the steady state, Section 3 the equilibrium conditions.

\item \texttt{results/solution.mat} --- everything \texttt{run.m} solved,
rewritten on each run.

\end{itemize}

\end{document}
